\documentclass[11pt]{article}
\usepackage[T1]{fontenc}
\usepackage{textcomp}
\usepackage[margin=0.75in]{geometry}
\usepackage{amsmath,amssymb,amsthm}
\usepackage{array,booktabs}
\usepackage{hyperref}
\setlength{\parindent}{0pt}
\newtheorem{theorem}{Theorem}
\theoremstyle{definition}
\newtheorem{definition}{Definition}
\title{Flow Proof Artifact: prop-24-parallelograms-equal-bases}
\author{Generated by \texttt{flow doc proof}}
\date{Flow Proof Book}
\begin{document}
\maketitle
Parallelograms on equal bases and in the same parallels are equal.\\[0.5em]
\textbf{Source.} euclid — Elements, Book I, Proposition 24\\[0.5em]
\bigskip
\noindent{\large \textbf{Derived fact 1.} \textit{Parallelograms on equal bases and in the same parallels are equal} \hfill the Euclidean plane \cdot Euclid Book I \cdot Proposition 24: parallelograms on equal bases and parallels are equal}\par
\smallskip\noindent\textit{Coordinate: the Euclidean plane \cdot Euclid Book I \cdot Proposition 24: parallelograms on equal bases and parallels are equal}\par
\smallskip\noindent\textit{Source: euclid — Elements, Book I, Proposition 24}\par
\smallskip\noindent\textit{Built on: proposition 4: side-angle-side congruence, for Book I of the Elements in on the Euclidean plane, proposition 14: lines with interior angles less than two right angles meet, for Book I of the Elements in on the Euclidean plane}\par
\medskip
\noindent\textbf{Goal.} Parallelograms on equal bases and in the same parallels are equal\par
\noindent$\displaystyle \text{parallelogram ABCD} = \text{parallelogram EFGH}$\par
\medskip
\noindent
\begin{tabular}{@{} >{\raggedright\arraybackslash}p{0.06\textwidth} >{\raggedright\arraybackslash}p{0.47\textwidth} >{\raggedleft\arraybackslash}p{0.41\textwidth} @{}}
\textbf{\#} & \textbf{Proof} & \textbf{Mathematics} \\
\midrule
\textbf{1.} & We prove that proposition 24: parallelograms on equal bases and parallels are equal for Book I of the Elements in the Euclidean plane. & \\[0.45em]
\textbf{2.} & Let diagonal\_AC = segment\_from\_A\_to\_C. & \\[0.45em]
\textbf{3.} & We invoke the derived fact: parallelogram\_ABCD\_and\_EFGH\_on\_equal\_bases\_BC\_and\_FG. & \\[0.45em]
\textbf{4.} & We invoke the axiom governing Book I of the Elements in the Euclidean plane: proposition 4: side-angle-side congruence, for Book I of the Elements in on the Euclidean plane. & \\[0.45em]
\textbf{5.} & We invoke the axiom governing Book I of the Elements in the Euclidean plane: proposition 14: lines with interior angles less than two right angles meet, for Book I of the Elements in on the Euclidean plane. & \\[0.45em]
\textbf{6.} & From step 2, step 3, step 4, and step 5, we can deduce that triangle ABC equals triangle DCB. & $\displaystyle \triangle ABC = \text{triangle DCB}$ \\[0.45em]
\textbf{7.} & We can deduce that parallelogram ABCD equals parallelogram EFGH. Hence proven. & $\displaystyle \text{parallelogram ABCD} = \text{parallelogram EFGH}$ \\[0.45em]
\bottomrule
\end{tabular}
\label{thm:Geometry-Euclid Book I-proposition_24_parallelograms_on_equal_bases_and_parallels_are_equal}
\par\medskip\hrule
\end{document}
